% SPDX-License-Identifier: Apache-2.0
\section{How a YubiKey works}
\label{app:yubikey}

This appendix explains the parts of a YubiKey that the study depends on.
It covers the YubiKey 5 series. Older models differ in details.

\subsection{One device, several applications}

A YubiKey holds several independent applications. Each one has its own keys,
its own PIN or password, and its own reset. Resetting the FIDO application
does not touch the PIV keys, and the reverse. Each application is reached
through one of the three USB interfaces in Table~\ref{tab:interfaces}.
Table~\ref{tab:apps} lists the applications.

\begin{table}[t]
\centering
\caption{YubiKey 5 applications}
\label{tab:apps}
\scriptsize
\begin{tabular}{@{}>{\raggedright\arraybackslash}p{13mm}>{\raggedright\arraybackslash}p{9mm}>{\raggedright\arraybackslash}p{24mm}>{\raggedright\arraybackslash}p{24mm}@{}}
\toprule
Application & USB interface & What it stores & How it proves \\
\midrule
FIDO U2F and FIDO2 & FIDO HID & a master key; discoverable credentials & signs a challenge after a touch \\
Yubico OTP & OTP HID (keyboard) & an AES-128 key and counters & types a 44-character one-time password \\
PIV & CCID & up to 24 key slots with certificates & signs or decrypts after a PIN \\
OpenPGP & CCID & signature, encryption and authentication keys & signs, decrypts or authenticates after a PIN \\
OATH & CCID & TOTP and HOTP secrets & computes a code for an app on the host \\
\bottomrule
\end{tabular}
\end{table}

The keys are made inside the YubiKey, and the device has no command that
exports a private key. A host can only ask the key to use a private key, for
example to sign. This is the reason why approach~C cannot work: a virtual
authenticator on the remote host cannot borrow the real key's private keys.

\subsection{FIDO credentials}

A FIDO credential is a key pair for one site, the \term{relying party}.
WebAuthn defines two operations~\cite{webauthn3}. Figure~\ref{fig:webauthn}
shows them.

\begin{figure}[t]
\centering
\begin{tikzpicture}[font=\scriptsize]
  \node[sbox,minimum width=14mm] (rp) at (0,0) {Site};
  \node[sbox,minimum width=14mm] (br) at (2.7,0) {Browser};
  \node[mem,minimum width=14mm] (yk) at (5.9,0) {YubiKey};
  \foreach \n in {rp,br,yk} \draw[black!30] (\n.south) -- ++(0,-5.35);
  \draw[fl] (0,-0.75) -- node[above]{challenge, RP ID} (2.7,-0.75);
  \draw[fl] (2.7,-1.35) -- node[above]{makeCredential} (5.9,-1.35);
  \node[warnb,font=\tiny,minimum height=4mm,inner sep=1.5pt] at (5.9,-1.85) {touch};
  \draw[fl] (5.9,-2.35) -- node[below]{public key, credential ID} (2.7,-2.35);
  \draw[fl] (2.7,-2.95) -- node[above]{signed response} (0,-2.95);
  \draw[densely dotted] (-0.6,-3.3) -- (6.6,-3.3);
  \node[anchor=west,font=\tiny] at (-0.6,-3.45) {later: login};
  \draw[fl] (2.7,-4.0) -- node[above]{getAssertion} (5.9,-4.0);
  \node[warnb,font=\tiny,minimum height=4mm,inner sep=1.5pt] at (5.9,-4.5) {touch};
  \draw[fl] (5.9,-5.0) -- node[below]{signature} (2.7,-5.0);
\end{tikzpicture}
\caption{Registration (top) and login (bottom). The browser sends the
challenge and the RP ID to the key; the key signs only after a touch.}
\label{fig:webauthn}
\end{figure}

\paragraph{Registration} The browser sends \code{makeCredential} with the
site's RP ID, a hash of the client data (which holds the site's challenge
and origin), and the user. After a touch, the key makes a new key pair. It
returns the public key and a \term{credential ID}, signed with its
attestation key.

\paragraph{Login} The browser sends \code{getAssertion} with the RP ID, a
new client data hash, and the credential IDs that the site knows. After a
touch, the key signs the authenticator data and the client data hash with
the credential's private key. The authenticator data holds a hash of the RP
ID, flags for user presence and user verification, and a signature counter.

\paragraph{Where the private key lives} A YubiKey makes a new key pair for
each registration. For a \term{non-discoverable} credential, it encrypts the
private key and some metadata with a master key that is unique to that
YubiKey and never leaves it. The result is the credential ID, which the site
stores~\cite{yk-keygen} (verified, firmware 4.4 and later). At login, the
site sends the credential ID back, and only the same YubiKey can decrypt it.
Such credentials therefore use no storage on the key. A
\term{discoverable} credential, also called a resident key or a passkey, is
stored on the key with the site's user handle, so a login needs no user
name~\cite{yk-resident} (verified). The key has room for a limited number of
these.

\paragraph{Presence and verification} A touch proves \term{user presence}: a
person is at the key. A PIN proves \term{user verification}: the right person
is at the key. The PIN is checked by the key itself, not by the host. Both
flags are part of the signed data, so the site can check them.

\paragraph{Why a phishing site gets no usable signature} The signed
data contains a hash of the RP ID, and the client data contains the origin.
A phishing site with a different name gets a signature that the real site
rejects. The browser chooses the RP ID. That is the reason why a relay that
trusts a compromised remote browser cannot stop phishing
(Section~\ref{sec:security}).

\subsection{Yubico OTP}

The OTP application acts as a USB keyboard. A touch makes it type a
44-character password~\cite{yk-otp} (verified). The first 12 characters are
the key's public ID, and they never change. The other 32 characters are
fields encrypted with the key's AES-128 key. The fields include a usage
counter and a session counter, so a server rejects a password that was used
before. The characters come from a small alphabet, called modhex, that types
the same on most keyboard layouts. The OTP application is not a goal of
this study.

\subsection{PIV}

PIV is a smart card standard from the US government. A YubiKey 4 or 5 has 24
PIV slots, and each holds a private key and an X.509
certificate~\cite{yk-piv} (verified). Four slots have defined uses:
\begin{itemize}
\item 9a, authentication: for example system login. One PIN entry permits
      many operations.
\item 9c, digital signature: the PIN is needed before every signature.
\item 9d, key management: decryption, for example of e-mail.
\item 9e, card authentication: for example physical access.
\end{itemize}
The PIN policy can be changed, and a slot can also ask for a touch. Hosts
reach PIV through the CCID interface, with PC/SC and a PKCS\#11 module. This
is the path that p11-kit forwards in approach~D.

\subsection{OpenPGP}

The OpenPGP application is an OpenPGP smart card~\cite{yk-pgp}. It holds up
to three keys: one to sign, one to decrypt, and one to authenticate. GnuPG
uses it through \code{scdaemon} and the CCID interface. A forwarded
gpg-agent lets a remote host ask for a signature, while the PIN entry and the
touch stay on the Chromebook (approach~D).

\subsection{What this means for forwarding}

Each application sits on one USB interface. A relay that forwards one
interface forwards only the applications behind it. The recommended relay
forwards the FIDO interface, so a compromised remote host can use FIDO
credentials while the relay runs, and nothing else. USB/IP forwards all
three interfaces, so the remote host can also use, or reset, the PIV and
OpenPGP applications and change the OTP configuration.
